\sscp{Words in Papuan languages similar to **mantəR ‘cuscus’}{14-04-04}
A handful of Papuan languages of the Bird’s Head Peninsula of
Papua and neighbouring islands have words formally similar to **mantəR which designate marsupials.
These languages are shown in \frf{14-05} with coloured squares,
with the words and glosses given in \trf{14-12}.
See appendix \ref{ch:SurMar} for the sources of these forms.

\begin{table}[h]\centering\tcp{Marsupial terms in Papuan languages resembling **mantəR}{14-12}
%\begin{adjustbox}{width=\textwidth}
	\begin{threeparttable}\st{6pt}
	\begin{tabular}{l|llll}\lsptoprule
Family	&	Language	&	ISO	&	form	&	gloss	\\	\midrule
\tsc{East Birds Head}	&	\ili{Moskona}	&	mtj	&	mioda	&	cuscus	\\
\tsc{East Birds Head}	&	\ili{Sougb}	&	kgr	&	meːsur	&	marsupial	\\
\tsc{\ili{Konda}-Yahadian}	&	\ili{Konda}	&	knd	&	‹märet›	&	tree wallaby	\\
\tsc{South Birds Head}	&	\ili{Kemberano}	&	bzp	&	metwetu	&	cuscus	\\
\tsc{South Birds Head}	&	\ili{Puragi}	&	pru	&	kamaʤi	&	tree wallaby	\\
\tsc{Greater West Bomberai}	&	\ili{Mbaham}	&	bdw	&	matabes\su{†}	&	cuscus	\\
\tsc{\ili{Yawa}-Saweru}	&	\ili{Yawa}	&	yva	&	maer	&	cuscus	\\
Isolate	&	\ili{Abun}	&	mnx	&	mit	&	kangaroo	\\		\lspbottomrule
	\end{tabular}
		\begin{tablenotes}\tnsize
			\item[†]	{Compare \ili{Mbaham} \it{matabes} with \ili{Erokwanas} \it{matafeis} ‘cuscus’.}
		\end{tablenotes}
	\end{threeparttable}
%\end{adjustbox}
\end{table}

In addition to the words in \trf{14-12} which denote marsupials,
as pointed out by \citet[269f]{Schapper2011},
there are formally similar words in Papuan languages of Wallacea
which refer to bats,
not cuscus, both of which are edible animals found in tree tops.
These languages are shown in \trf{14-13}.
At this point we simply note the remarkable similarities in form
to **mantəR ‘cuscus’ in unrelated Papuan languages,
rather than asserting a connection.
See more details in appendix \ref{ch:SurMar}.
%Forms without a citation in \trf{14-13} are from \citet{UsherSchapper2022}.


\begin{table}[h]\centering\tcp{\tsc{Greater West Bomberai} words similar to **mantəR}{14-13}
%\begin{adjustbox}{width=\textwidth}
	\begin{threeparttable}
		%\begin{adjustbox}{width=\textwidth}
	%\st{2pt}\begin{tabular}{\textwidth}{l|lll}\lsptoprule
	\st{2.75pt}\begin{tabularx}{\textwidth}{l|ll>{\raggedright\arraybackslash}X}\lsptoprule
Language	&	Form	&	local gl.	&	Source	\\	\midrule
P.-G. W. B.	&	*mandəl	&	flying fox	&	\citet[493]{UsherSchapper2022}	\\	\midrule
\ili{Iha}	&	manda	&	bat	&	\citet[122]{SmitsVoorhoeve1998}	\\	\midrule
P.-TAP\su{†}	&	*madel	&	flying fox	&	\citet[493]{UsherSchapper2022}	\\
P.-E. Timor	&	*mada	&	flying fox	&	\citet[493]{UsherSchapper2022}	\\	\midrule
\ili{Fataluku}	&	mat͡sa	&	bat	&	\citet[189]{Heston2015}	\\
\ili{Oirata}	&	maʈa-mailur	&	large bat	&	\citet{DeJosselinDeJong1937} cited in \citet[493]{UsherSchapper2022}	\\	\midrule
P.-AP	&	*madəl	&	bat	&	\citet[78]{HoltonRobinson2017}	\\	\midrule
\ili{Wersing}	&	məduku	&	bat	&	Johnny Banamtuan pers. comm. Oct 2020	\\
\ili{Sawila}\su{†}	&	madiiku 	&	bat	&	Johnny Banamtuan pers. comm. Oct 2020	\\
\ili{Kamang}	&	maˈtei	&	bat	&	\citet[87]{SchapperManimau2011}	\\
\ili{Abui}	&	maˈrɛl	&	bat	&	Dialect variants: \it{mate, matii}	\\
\ili{Kafoa}	&	marel	&	bat	&	\citet[85]{KratochvilDelpada2008}	\\
\ili{Klon}	&	mədel	&	bat	&	\citet[269]{Schapper2011}; Johnny Banamtuan pers. comm. Oct 2020	\\
\ili{Teiwa}	&	medi	&	bat	&	\citet[35]{KlamerSir2011}	\\
\ili{Klamu}\su{‡}	&	[ˈmărˑɑ]	&	bat	&	Grimes fieldnotes; see \srf{02-04-06} 	\\
\ili{Papuna}	&	marie	&	bat	&	\citet{Delpada2011}	\\
\ili{Mauta} {\&}	&	madːe	&	bat	&	\citet[92]{HoltonKoly2008} `k.o. slightly	\\
\ili{Tubbe}	&	madːe	&	bat	&	furry bat'	\\
\ili{Lamma}	&	madːal	&	flying fox	&	\citet[92]{HoltonKoly2008} Indonesian \it{kalung}	\\
\lspbottomrule
	\end{tabularx}
%	\end{tabular}
		\begin{tablenotes}\tnsize
			\item[†]	{\citet[269]{Schapper2011} provided a Proto-TAP reconstruction *mader ‘bat’.
								The final consonant of the reconstruction is based on \ili{Sawila} \it{madiirang} ‘bat’.
								But that form is incorrect.
								\ili{Sawila} \it{madiirang} means ‘sharp’.
								One item up in the word list is \it{madiiku} ‘bat, (Indonesian \it{kelelawar})’.}
			\item[‡]	{\citet[493]{UsherSchapper2022} list Nedebang (= \ili{Klamu}) \it{marːa} ‘bat’.}
								%Nedebang is an alternate name for \ili{Klamu}.}
		\end{tablenotes}
	\end{threeparttable}
%\end{adjustbox}
\end{table}

We note the similarity of forms across Austronesian
and Papuan language families,
both referring to edible mammals dwelling in treetops.
We specifically note the similarities of the Proto-Timor-\ili{Alor}-\ili{Pantar}
*madel ‘flying fox’, as being remarkably close in form to \tsc{Timor-Babar} **mantəR/**məntəR
‘cuscus’ and \tsc{\ili{Central} Timor} **mantəR ‘cuscus’,
accounting for nearby Austronesian languages with which they
are in regular contact.

\citet[486]{UsherSchapper2022} see the Papuan Proto-Greater West
Bomberai word *mandəl ‘flying fox’ (see \frf{14-06}) as being
the ultimate loan source into some nearby Austronesian languages,
including: \ili{Kowiai} \it{mada-mada} ‘small bat’,
\ili{Sekar} \it{mada} ‘fruit bat’,
Patipi (\ili{Onin}) \it{mad-mada} ‘flying fox’,
\ili{Uruangnirin} \it{mada} ‘bat’ (all in the Bomberai region),
along with Waima{\Q}a \it{mada} ‘bat’
and \ili{Naueti} \it{mada} ‘bat’ in the Timor region.
Once again, contact is a factor in the spread of these related terms.
It very likely happened independently in the two regions,
and there is no need to reconstruct a proto-form in a phantom
parent language to account for the similarities.
We also note that without an awareness of the Papuan forms,
one might easily make such a reconstruction.

\citet[269f]{Schapper2011} noted that \ili{Bati} \it{manuk faniki} ‘cuscus’
(< *manuk ‘chicken, bird’; *paniki ‘fruit bat,
flying fox’) attests a link between the word for `cuscus'
and flying creatures, including bats.\footnote{
		The related cognate term in \ili{Buru} \it{manu-t}
		includes birds, bats, and butterflies.
		That is, it is a generic term covering winged creatures whose
		wing motion is slow enough to observe.
		\ili{Buru} \it{man grihit} ‘flying fox’
		and \it{man koi} ‘small cave bat’ show regular truncation of
		the generic term also found for birds,
		such as \it{man keho} ‘mallee fowl’,
		\it{man kohot} ‘crow’, \it{man kumul} ‘dove,
		pigeon’, and \it{man lapa} ‘butterfly’, etc.
		See \citet{GrimesGrimes2020}.}
In addition to \ili{Bati}, \ili{Masiwang} \it{mani} ‘cuscus’
is also from PMP *manuk `chicken'
(ultimate *u > \it{i} is regular as shown in Tables
\ref{tab:11-11} and \ref{tab:11-12})

At this point we are merely observing the remarkable similarities
in form, and possible links in meaning,
which are collectively too close to simply dismiss out of hand,
and are also evidence of contact in this semantic domain
and with these terms.

\begin{figure}[h]\centering
	\includegraphics[width=0.6\columnwidth]{14-06.jpg}
	\fcp{Flying fox, fruit bat}{14-06}
\end{figure}